We compare three capacity policies over a shared evidence ordering: fixed-$k$
uses a fixed evidence count, minmax$(k_{\min},k_{\max})$ allows
structure-conditioned stopping within a bounded range, and budget-$B$ limits
the verifier-visible context by tokens.

\begin{figure}[t]
  \centering
  \includegraphics[width=\columnwidth]{evidence_capacity_sensitivity.pdf}
  \caption{Evidence-capacity sensitivity on the LIAR-RAW test set
  ($n{=}1{,}251$). Points show Macro-F1 against the mean prompt tokens presented
  to the verifier; higher Macro-F1 with fewer prompt tokens is preferable.
  Colors, markers, and lines distinguish fixed-$k$, minmax, and budget policies. Point labels
  identify the exact setting: F$k$ denotes fixed-$k$, M$a$--$b$ denotes
  minmax$(a,b)$, and B$n$ denotes token budget $n$. Lines are visual guides
  rather than fitted trends. The blue star marks the best observed
  configuration, minmax$(5,10)$.}
  \label{fig:capacity-sensitivity}
\end{figure}

Minmax$(5,10)$ provides the best observed operating point, reaching 36.66
Macro-F1 with 677.90 mean prompt tokens; longer prompts do not consistently
improve verification performance.
